\chapter{Evenredigheid en gegevens}\label{ch:g6:propdata}

Als drie schriften $6$ euro kosten, kosten zes schriften $12$: dubbel zoveel
schriften, dubbel zoveel geld. Grootheden die zich zo gedragen, zijn
\emph{\omterm{def:g6:propdata:def}{evenredig}} --- een van de nuttigste ideeën uit de hele wiskunde, dat in
\cref{ch:g7:prop} verder wordt uitgewerkt. Het hoofdstuk eindigt met het lezen
en tekenen van eenvoudige diagrammen.

\section{Evenredige grootheden}

\begin{definition}[Evenredigheid]\label{def:g6:propdata:def}
Twee grootheden zijn \emph{evenredig}\index{evenredigheid} als je de waarden van
de ene uit de waarden van de andere krijgt door steeds met \emph{hetzelfde} getal
te vermenigvuldigen; dat getal heet de \emph{evenredigheidsconstante}.
\end{definition}

\begin{example}\label{ex:g6:propdata:table}
Schriften van $2$ euro per stuk:
\begin{center}
\renewcommand{\arraystretch}{1.3}
\begin{tabular}{l|cccc}
schriften & $3$ & $5$ & $8$ & $12$ \\
\hline
prijs (euro) & $6$ & $10$ & $16$ & $24$ \\
\end{tabular}
\end{center}
Elke prijs is het aantal schriften $\times 2$: de constante is $2$ (de prijs per
stuk). Een snelle controle of een tabel \omterm{def:g6:propdata:def}{evenredig} is: alle ``kolomquotiënten''
$\frac{6}{3}, \frac{10}{5}, \frac{16}{8}, \frac{24}{12}$ moeten gelijk zijn ---
hier zijn ze alle $2$.
\end{example}

\begin{example}[Een tegenvoorbeeld]\label{ex:g6:propdata:nonexample}
Leeftijd en lengte zijn niet \omterm{def:g6:propdata:def}{evenredig}: een kind van $12$ is niet twee keer zo
lang als een kind van $6$. Controleer met getallen: $1.50$ m op $12$ jaar en
$1.15$ m op $6$ jaar geven \omterm{def:g3:division:remainder}{quotiënten} $\frac{1.50}{12} = 0.125$ en
$\frac{1.15}{6} \approx 0.19$ --- niet gelijk.
\end{example}

\begin{method}[Een evenredigheidstabel aanvullen]\label{met:g6:propdata:complete}
Drie manieren, die je vrij mag kiezen:
\begin{enumerate}
  \item \emph{de constante}: zoek de vermenigvuldiger uit één volledige kolom en
        gebruik hem bij de andere;
  \item \emph{kolommen}: een kolom mag met een getal vermenigvuldigd worden, of
        twee kolommen mogen opgeteld worden, om een nieuwe kolom te maken;
  \item \emph{terug naar één}: zoek eerst de waarde voor \emph{één} stuk en
        vermenigvuldig daarna.
\end{enumerate}
\end{method}

\begin{example}\label{ex:g6:propdata:methods}
Vijf dezelfde mokken kosten $15$ euro; wat kosten acht mokken?

\emph{Terug naar één}: één mok kost $15 \div 5 = 3$ euro, dus kosten acht mokken
$8 \times 3 = 24$ euro.

\emph{Kolommen}: $8 = 5 + 3$; drie mokken kosten $\frac{3}{5}$ van $15$, dus
$9$; acht mokken kosten dan $15 + 9 = 24$ euro. Hetzelfde antwoord, zoals het
hoort.
\end{example}

\section{Percentages}

\begin{definition}[Percentage]\label{def:g6:propdata:percent}
``$t\,\%$ van een hoeveelheid'' betekent de \omterm{def:g6:fractions:def}{breuk} $\frac{t}{100}$ ervan:
$25\,\%$ van $60$ is $\frac{25}{100} \times 60 = 15$. Een percentage is een
\omterm{def:g6:propdata:def}{evenredigheid} met constante $\frac{t}{100}$.
\end{definition}

\begin{example}\label{ex:g6:propdata:percent}
Een klas van $30$ leerlingen bestaat voor $40\,\%$ uit meisjes. Het aantal
meisjes, stap voor stap: $40\,\%$ betekent $\frac{40}{100} = 0.4$, en
$0.4 \times 30 = 12$ meisjes. Handige ijkpunten: $50\,\%$ is de \omterm{def:g3:division:half}{helft}, $25\,\%$
een \omterm{def:g3:division:half}{kwart}, $10\,\%$ een tiende --- dus is $10\,\%$ van $30$ gelijk aan $3$, en
$40\,\% = 4 \times 10\,\%$ geeft $4 \times 3 = 12$: hetzelfde antwoord, uit het
hoofd gerekend.
\end{example}

\section{Diagrammen lezen en tekenen}

\begin{method}[Een tabel of diagram lezen]\label{met:g6:propdata:read}
Wat het plaatje ook is --- tabel, staafdiagram, lijndiagram:
\begin{enumerate}
  \item lees eerst de titel en de eenheden op beide assen;
  \item zoek voor een vraag de juiste staaf of kolom, en lees de waarde
        zorgvuldig af tegen de streepjes;
  \item vergelijk staven op hoogte --- maar controleer of de as bij $0$ begint,
        anders kan het plaatje misleiden.
\end{enumerate}
\end{method}

\begin{example}\label{ex:g6:propdata:chart}
Het staafdiagram hieronder toont het aantal boeken dat in één week in de
schoolbibliotheek is uitgeleend.

\begin{omfigure}
\begin{tikzpicture}
\begin{axis}[
  width=8.6cm, height=5.2cm,
  ybar, bar width=16pt,
  symbolic x coords={ma,di,wo,do,vr},
  xtick=data,
  ymin=0, ymax=22,
  ytick={5,10,15,20},
  ylabel={uitgeleende boeken},
  tick label style={font=\small},
  label style={font=\small}]
  \addplot[fill=omDef!30, draw=omDef] coordinates
    {(ma,12) (di,8) (wo,17) (do,6) (vr,15)};
\end{axis}
\end{tikzpicture}

{\small Een staafdiagram: één staaf per dag, met een hoogte gelijk aan het
aantal.}
\end{omfigure}

Aflezen: de drukste dag is woensdag ($17$ boeken). Dinsdag en donderdag samen
zijn $8 + 6 = 14$ boeken --- minder dan vrijdag alleen ($15$). Totaal voor de
week: $12 + 8 + 17 + 6 + 15 = 58$ boeken.
\end{example}

\begin{omfigure}
\begin{tikzpicture}
\begin{axis}[omaxis,
  width=8.2cm, height=5.4cm,
  xmin=0, xmax=9.4, ymin=0, ymax=28,
  xlabel={schriften}, ylabel={prijs (euro)},
  xtick={1,2,3,4,5,6,7,8,9}, ytick={4,8,12,16,20,24}]
  \addplot[omDef, thick, domain=0:9] {3*x};
  \addplot[omDef, only marks, mark=*, mark size=1.5pt] coordinates
    {(1,3) (2,6) (3,9) (4,12) (5,15) (6,18) (7,21) (8,24)};
  \draw[dashed, omThm] (axis cs:6,0) -- (axis cs:6,18)
    -- (axis cs:0,18);
\end{axis}
\end{tikzpicture}

{\small Evenredige grootheden hebben een heel goed te herkennen grafiek: de
punten liggen op één rechte \omterm{def:g6:lines:objects}{lijn} door de oorsprong (hier schriften van $3$ euro
per stuk; de stippellijnen lezen de prijs van $6$ schriften af: $18$ euro).}
\end{omfigure}

\section{Oefeningen}

\begin{exercise}[$\star$]\label{exo:g6:propdata:1}
Is de tabel \omterm{def:g6:propdata:def}{evenredig}? Verantwoord het door \omterm{def:g3:division:remainder}{quotiënten} uit te rekenen.
\begin{center}
\renewcommand{\arraystretch}{1.2}
\begin{tabular}{l|ccc}
kg appels & $2$ & $3$ & $5$ \\
\hline
prijs (euro) & $5$ & $7.5$ & $12.5$ \\
\end{tabular}
\qquad
\begin{tabular}{l|ccc}
leeftijd (jaar) & $2$ & $4$ & $8$ \\
\hline
lengte (cm) & $85$ & $103$ & $130$ \\
\end{tabular}
\end{center}
\end{exercise}

\begin{exercise}[$\star$]\label{exo:g6:propdata:2}
Vier croissants kosten $4.80$ euro. Zoek de prijs van één croissant, en daarna
die van zeven croissants.
\end{exercise}

\begin{exercise}[$\star$]\label{exo:g6:propdata:3}
Vul de evenredigheidstabel aan (eerst de constante!):
\begin{center}
\renewcommand{\arraystretch}{1.2}
\begin{tabular}{l|cccc}
liter brandstof & $10$ & $25$ & $?$ & $60$ \\
\hline
prijs (euro) & $18$ & $?$ & $72$ & $?$ \\
\end{tabular}
\end{center}
\end{exercise}

\begin{exercise}[$\star$]\label{exo:g6:propdata:4}
Een auto verbruikt $6$ L brandstof per $100$ km. Hoeveel brandstof voor $250$
km? En voor $350$ km? Hoe ver kan hij komen met $27$ L?
\end{exercise}

\begin{exercise}[$\star$]\label{exo:g6:propdata:5}
Reken uit het hoofd, met de ijkpunten van \cref{ex:g6:propdata:percent}:
$50\,\%$ van $84$; \quad $25\,\%$ van $200$; \quad $10\,\%$ van $63$; \quad
$30\,\%$ van $70$.
\end{exercise}

\begin{exercise}[$\star$]\label{exo:g6:propdata:6}
Op een school met $450$ leerlingen eet $60\,\%$ in de kantine. Hoeveel
leerlingen zijn dat? En hoeveel niet?
\end{exercise}

\begin{exercise}[$\star$]\label{exo:g6:propdata:7}
Gebruik het staafdiagram van \cref{ex:g6:propdata:chart}: op welke dagen werden
er minder dan $10$ boeken uitgeleend? Hoeveel boeken meer werden er op woensdag
uitgeleend dan op donderdag?
\end{exercise}

\begin{exercise}[$\star$]\label{exo:g6:propdata:8}
De temperaturen om twaalf uur 's middags waren van maandag tot vrijdag $14$,
$16$, $13$, $17$ en $20$ graden. Teken een staafdiagram van deze gegevens (kies
zinnige streepjes) en lees de warmste dag af.
\end{exercise}

\begin{exercise}[$\star\star$]\label{exo:g6:propdata:9}
Een recept voor $6$ personen vraagt $450$ g bloem en $3$ eieren. Pas het aan
voor $10$ personen. (Denk bij de eieren na voordat je een decimaal aantal eieren
opschrijft!)
\end{exercise}

\begin{exercise}[$\star\star$]\label{exo:g6:propdata:10}
Bij een vaste snelheid rijdt een wielrenster $24$ km in $1$ uur.
\begin{enumerate}
  \item Hoe ver komt ze in $2$ u? En in een half uur? En in $1$ u $30$ min?
  \item Hoe lang heeft ze nodig voor $60$ km?
\end{enumerate}
\end{exercise}

\begin{exercise}[$\star\star$]\label{exo:g6:propdata:11}
Een winkel biedt ``$20\,\%$ korting op alles''. Reken de korting en de nieuwe
prijs uit voor: een bal van $15$ euro; een racket van $40$ euro. Is de nieuwe
prijs \omterm{def:g6:propdata:def}{evenredig} met de oude prijs? Wat is de constante?
\end{exercise}

\begin{exercise}[$\star\star\star$]\label{exo:g6:propdata:12}
Twee kaarsen van dezelfde hoogte worden op hetzelfde moment aangestoken. De
dikke brandt in $6$ uur volledig op, de dunne in $4$ uur, elk met haar eigen
vaste snelheid. Na hoeveel tijd is de dunne kaars precies half zo hoog als de
dikke? (Schrijf de resterende hoogtes na $t$ uur als \omterm{def:g6:fractions:def}{breuken} van de
beginhoogte.)
\end{exercise}

\section{Opgave: kaarten, plattegronden en hoe je liegt met een diagram}

\begin{problem}[{Weekendopgave --- de schaal van een kaart is een
evenredigheid: lengtes volgen de constante, oppervlaktes haar kwadraat, en
slecht getekende diagrammen bedriegen het oog}]
\label{pb:g6:propdata:1}
Elke kaart, plattegrond en maquette volgt één regel: echte lengtes en getekende
lengtes zijn \emph{\omterm{def:g6:propdata:def}{evenredig}} (\cref{def:g6:propdata:def}). De constante heeft
een beroemde naam --- de \emph{schaal}\index{schaal} --- en er zitten twee
valstrikken in verstopt die deze opgave met opzet laat toeslaan: \omterm{def:g6:measure:area}{oppervlaktes}
volgen de schaal \emph{niet}, en percentages hangen, net als diagrammen,
volledig af van waartegen ze gemeten worden.

\medskip
\noindent\textbf{Deel I --- Een kaart lezen.}
Een wandelkaart heeft schaal $1:100\,000$: elke centimeter op de kaart staat
voor $100\,000$ centimeter op de grond.
\begin{enumerate}
  \item Zet $100\,000$ cm om in kilometers. Waar staat $1$ cm van deze kaart
        voor, in km?
  \item Twee dorpen liggen op de kaart $7.5$ cm van elkaar; de oever van een
        meer is een kromme van $12$ cm. Geef de echte afstanden.
  \item Een kaarsrechte Romeinse weg is $20$ km lang. Hoe lang is hij op de
        kaart?
  \item Een tweede kaart kondigt ``$1$ cm voor $5$ km'' aan. Schrijf haar schaal
        in de vorm $1: \,?$. Welke van de twee kaarten toont meer detail voor
        hetzelfde gebied?
  \item Maak een klein tabelletje (kaartafstand $1$, $2$, $7.5$ cm tegen echte
        afstand) voor de wandelkaart, en leg uit waarom een schaal precies een
        \omterm{def:g6:propdata:def}{evenredigheid} is in de zin van \cref{def:g6:propdata:def}. Wat is de
        constante die centimeters-op-de-kaart in kilometers omzet?
\end{enumerate}

\medskip
\noindent\textbf{Deel II --- De plattegrond van Zoë's slaapkamer.}
Zoë tekent een plattegrond van haar slaapkamer --- een rechthoek van $4$ m bij
$3$ m --- op schaal $1:50$.
\begin{enumerate}[resume]
  \item Wat zijn de afmetingen van de kamer op de plattegrond?
  \item Haar bed meet $2$ m bij $0.9$ m, en de deuropening is $80$ cm breed.
        Geef alle drie de maten op de plattegrond.
  \item Nu de valstrik. Reken de echte \omterm{def:g6:measure:area}{oppervlakte} van de kamer in m$^2$ uit, en
        daarna de \omterm{def:g6:measure:area}{oppervlakte} van de plattegrond in cm$^2$. Zet de echte
        \omterm{def:g6:measure:area}{oppervlakte} om in cm$^2$ (\cref{def:g6:measure:area}) en deel door de
        \omterm{def:g6:measure:area}{oppervlakte} van de plattegrond: is het \omterm{def:g3:division:remainder}{quotiënt} $50$?
  \item Verklaar het getal dat je vond: op een plattegrond $1:50$ staat elke
        vierkante centimeter papier voor een echt vierkant van $50$ cm bij $50$
        cm. Hoeveel echte cm$^2$ is dat? Formuleer de regel: als lengtes door
        $50$ gedeeld worden, worden \omterm{def:g6:measure:area}{oppervlaktes} gedeeld door\,\dots
  \item Een rond vloerkleed bedekt $6$ cm$^2$ op de plattegrond. Welke echte
        \omterm{def:g6:measure:area}{oppervlakte} bedekt het, in cm$^2$ en daarna in m$^2$?
\end{enumerate}

\medskip
\noindent\textbf{Deel III --- Hoe je liegt met een diagram (en met een
percentage).}
\begin{enumerate}[resume]
  \item Een winkel verkocht op zaterdag voor $105$ euro en op zondag voor $110$.
        De eigenaar tekent een staafdiagram waarvan de verticale as bij $100$
        begint: de staaf van zaterdag stijgt $5$ kleine eenheden, die van zondag
        $10$. Wat suggereert het plaatje over zondag, en wat is de waarheid?
        (Reken de stijging van zondag uit als percentage van zaterdag, op het
        hele procent.) Welke waarschuwing van \cref{met:g6:propdata:read} heeft
        de eigenaar genegeerd?
  \item Teken (of beschrijf) het eerlijke diagram, met de as beginnend bij $0$:
        hoe verhouden de twee staven zich nu?
  \item Een verzameling groeit van $40$ naar $60$ postzegels: reken de groei uit
        als percentage van de beginwaarde. Daarna krimpt ze weer van $60$ naar
        $40$: reken de daling uit als percentage van \emph{haar} beginwaarde.
        Waarom zijn de twee antwoorden verschillend, voor dezelfde $20$
        postzegels?
  \item Uitverkoop: een jas van $100$ euro krijgt ``$20\,\%$ korting'', en aan de
        kassa nog eens ``$10\,\%$ korting op de verlaagde prijs''. Reken de
        eindprijs stap voor stap uit. Is de totale korting $30\,\%$? Leg de
        koper uit wat het echt is.
  \item Finale, terug op de kaart van vraag 4 ($1$ cm voor $5$ km): een bos beslaat
        $3$ cm$^2$ op die kaart. Wat is zijn echte \omterm{def:g6:measure:area}{oppervlakte}, in km$^2$? Besluit
        met de regel van deze hele opgave: lengtes volgen de schaal,
        \omterm{def:g6:measure:area}{oppervlaktes} haar\,\dots
\end{enumerate}
\end{problem}
